% This file was converted to LaTeX by Writer2LaTeX ver. 1.9.9
% see http://writer2latex.sourceforge.net for more info
\documentclass{article}
\usepackage{calc,amsmath,amssymb,amsfonts}
\usepackage[T1]{fontenc}
\usepackage[italian]{babel}
\usepackage[style=numeric,backend=biber]{biblatex}
\author{HP}
\date{2026-07-26}
\begin{document}
\section{Prefazione}
Ogni libro nasce da una domanda.

Alcuni cercano di spiegare una tecnologia. Altri raccontano una storia. Altri ancora desiderano tramandare
un'esperienza.

Questo libro nasce da tutte e tre queste esigenze.

Parleremo di un computer che ha segnato un'epoca. Impareremo due linguaggi di programmazione che hanno formato
generazioni di sviluppatori. Costruiremo, passo dopo passo, un videogioco completo. Ma, soprattutto, proveremo a
ritrovare qualcosa che il tempo non è riuscito a cancellare: la curiosità.

Negli anni Ottanta, programmare un computer non era soltanto un passatempo. Era un modo di esplorare un territorio
sconosciuto. Ogni comando scoperto sembrava aprire una porta. Ogni errore costringeva a ragionare. Ogni piccolo
successo regalava una soddisfazione difficile da descrivere a chi non l'ha vissuta.

Il Commodore 64 è stato, per migliaia di ragazzi, il primo laboratorio di informatica. Non chiedeva di essere
semplicemente utilizzato. Invitava a fare domande. Spingeva a smontare le certezze, a sperimentare, a sbagliare e a
riprovare. In un'epoca in cui Internet non esisteva e le informazioni si cercavano sulle riviste, nei manuali o
attraverso il passaparola tra amici, ogni nuova scoperta aveva il sapore di una conquista personale.

Molti di quei ragazzi sono oggi professionisti dell'informatica. Altri hanno seguito strade completamente diverse.
Eppure, quando rivedono quella tastiera color crema o la schermata blu con il cursore lampeggiante, succede qualcosa di
particolare. Per un istante il tempo sembra fare un passo indietro.

Questo libro è dedicato anche a quel momento.

Non vogliamo limitarci a spiegare come funziona il Commodore 64. Vorremmo accompagnarti nel percorso che portava un
ragazzo qualunque a chiedersi una domanda apparentemente semplice e, proprio per questo, straordinaria: come si
costruisce un videogioco?

Questa domanda guiderà ogni pagina che seguirà.

Non inizieremo studiando teoria per poi cercare un'applicazione pratica. Faremo esattamente il contrario. Avremo un
obiettivo concreto: realizzare un videogioco. Ogni volta che, durante lo sviluppo, incontreremo un nuovo problema, ci
fermeremo per comprenderlo. Se avremo bisogno degli sprite, scopriremo come funzionano gli sprite. Se il BASIC non sarà
più sufficiente, entreremo nel mondo dell'Assembly. Se vorremo rendere il gioco più veloce o più spettacolare,
esploreremo l'hardware del Commodore 64 fino a comprenderne i meccanismi più affascinanti.

In altre parole, non studieremo per programmare un videogioco.

Programmeremo un videogioco per imparare.

È una differenza sottile, ma cambia completamente il modo di affrontare questo percorso.

Anche il ruolo del secondo autore merita qualche parola.

La presenza di ChatGPT potrebbe sorprendere qualcuno. È naturale chiedersi quale sia il contributo di un'intelligenza
artificiale in un libro dedicato a un computer del 1982.

La risposta è semplice.

Questo libro non è stato scritto da una persona che ha delegato il proprio lavoro a un'intelligenza artificiale. È nato
da un dialogo continuo. Da domande, dubbi, verifiche, confronti, riscritture e idee migliorate insieme. Abbiamo scelto
di dichiararlo apertamente perché crediamo che questo rappresenti un nuovo modo di costruire conoscenza: collaborare
con strumenti moderni senza rinunciare al pensiero critico, alla creatività e alla responsabilità delle proprie scelte.

Nel corso della lettura, ChatGPT non assumerà mai il ruolo di chi possiede tutte le risposte. Preferirà comportarsi come
il compagno di laboratorio che tutti avremmo voluto avere: quello che suggerisce una strada, propone un esperimento,
invita a osservare un dettaglio, ma lascia sempre a noi il piacere della scoperta.

Per questo motivo il vero protagonista del libro non è il Commodore 64.

Non è ChatGPT.

E non siamo nemmeno noi che l'abbiamo scritto.

Il protagonista sei tu.

Se possiedi ancora il tuo Commodore 64, proveremo ad aiutarlo a tornare in vita con tutte le precauzioni necessarie.

Se preferirai utilizzare un emulatore, costruiremo insieme un laboratorio moderno capace di offrirti strumenti che negli
anni Ottanta avremmo considerato fantascienza.

Se, invece, non hai mai visto un Commodore 64 acceso, speriamo che queste pagine riescano a trasmetterti almeno una
parte dell'entusiasmo che quella macchina seppe accendere in un'intera generazione.

Alla fine del viaggio avremo costruito un videogioco.

Ma non sarà questo il risultato più importante.

Il vero traguardo sarà aver imparato a guardare un computer con gli occhi della curiosità, ricordando che dietro ogni
programma, ogni pixel e ogni suono c'è sempre un'idea nata nella mente di qualcuno.

È da quella curiosità che nascono i programmatori.

Ed è proprio da lì che inizierà anche il nostro viaggio.


\bigskip
\end{document}
